\documentclass[11pt]{article}

% Plantilla maestra para guías del instructor
\input{~/Dropbox/Latex/TemplateGuide/template-guide.tex}

\makeatletter
\renewcommand*\l@subsection{\@dottedtocline{2}{1.5em}{3.8em}}
\makeatother

\renewcommand{\CourseName}{Diseño Web 1}
\renewcommand{\CourseCode}{ADSO-CTMA}
\renewcommand{\Faculty}{Análisis y Desarrollo de Software}
\renewcommand{\Institution}{Servicio Nacional de Aprendizaje SENA}
\renewcommand{\Instructor}{Deimer Miranda Montoya, MSc.(c).}
\renewcommand{\Email}{damiranda@sena.edu.co}
\renewcommand{\GuideTitle}{Guía 4: Diseño responsivo e interacción visual}
\renewcommand{\ClassName}{Responsive, estados y animaciones}
\renewcommand{\AcademicPeriod}{2026-3T}
\renewcommand{\InstitutionLogo}{Template/LogoSena.png}

\begin{document}
\thispagestyle{firstpage}
\GuideHeader
\GuideInfo{Clases 9 y 10}{Presencial con explicación, práctica guiada y trabajo autónomo}{Interfaz adaptable con estados, transiciones y movimiento accesible}

\section{Presentación}
Esta guía parte de una interfaz que ya tiene estructura y layout. El nuevo objetivo es hacer que conserve contenido, legibilidad y control en diferentes tamaños de pantalla y, después, agregar estados, transiciones y animaciones con una función comunicativa clara.

\section*{Resultados de aprendizaje de la guía}
Al finalizar, el aprendiz podrá:
\begin{itemize}
    \item explicar el propósito del viewport en dispositivos móviles;
    \item combinar medidas fluidas, límites y recursos adaptables;
    \item escribir media queries a partir de necesidades reales del contenido;
    \item aplicar un enfoque mobile first;
    \item diseñar estados visibles para interacción por puntero y teclado;
    \item utilizar transformaciones y transiciones con moderación;
    \item construir animaciones mediante \texttt{@keyframes};
    \item respetar \texttt{prefers-reduced-motion} cuando el movimiento no sea esencial.
\end{itemize}

\section*{Conocimientos previos}
Se espera que el aprendiz pueda construir una página con HTML semántico y organizarla mediante propiedades CSS, Flexbox o Grid.

\Warning[Alcance de esta guía]{
No se introduce JavaScript. La interacción trabajada corresponde a estados y movimiento que CSS puede expresar por sí mismo.
}

\tableofcontents
\newpage

\section{De una pantalla fija a una interfaz adaptable}
\begin{center}
\begin{tikzpicture}[
    pantalla/.style={draw=BrandBlue, very thick, rounded corners, fill=white},
    bloque/.style={fill=PracticeBody, draw=PracticeTitle, thick}
]
    \node[font=\bfseries] at (1.8,4.5) {Móvil};
    \draw[pantalla] (0,0) rectangle (3.6,4.1);
    \draw[bloque] (0.35,3.25) rectangle (3.25,3.7);
    \draw[bloque] (0.35,2.15) rectangle (3.25,3.0);
    \draw[bloque] (0.35,1.05) rectangle (3.25,1.9);
    \draw[bloque] (0.35,0.3) rectangle (3.25,0.8);

    \node[font=\bfseries] at (8.1,4.5) {Pantalla amplia};
    \draw[pantalla] (4.5,0) rectangle (11.7,4.1);
    \draw[bloque] (4.85,3.25) rectangle (11.35,3.7);
    \draw[bloque] (4.85,1.15) rectangle (6.8,3.0);
    \draw[bloque] (7.15,1.15) rectangle (9.1,3.0);
    \draw[bloque] (9.45,1.15) rectangle (11.35,3.0);
    \draw[bloque] (4.85,0.3) rectangle (11.35,0.8);
\end{tikzpicture}
\end{center}

\Idea[La pregunta correcta]{
Diseño responsive no significa ``hacer una versión para celular''. Significa permitir que el contenido reorganice su presentación cuando el espacio disponible cambia.
}
\section{Laboratorio secuencial: Ruta Verde}

Se desarrollará una página informativa sobre recorridos ecológicos. La estructura
se mantendrá durante toda la guía y el CSS evolucionará de una columna a una
composición responsiva con movimiento moderado.

\subsection{Archivos}

\begin{consolecode}
ruta-verde/
├── index.html
├── css/
│   └── styles.css
└── img/
    ├── portada.jpg
    ├── bosque.jpg
    ├── cascada.jpg
    └── mirador.jpg
\end{consolecode}

\begin{htmlfile}{index.html}
<!DOCTYPE html>
<html lang="es">
<head>
    <meta charset="UTF-8">
    <meta name="viewport" content="width=device-width, initial-scale=1.0">
    <title>Ruta Verde</title>
    <link rel="stylesheet" href="css/styles.css">
</head>
<body>
    <header class="encabezado">
        <a class="marca" href="#inicio">Ruta Verde</a>
        <nav aria-label="Navegación principal">
            <ul class="menu">
                <li><a href="#rutas">Rutas</a></li>
                <li><a href="#recomendaciones">Recomendaciones</a></li>
                <li><a href="#contacto">Contacto</a></li>
            </ul>
        </nav>
    </header>

    <main id="inicio">
        <section class="hero">
            <div class="hero__contenido">
                <p class="etiqueta">Exploración responsable</p>
                <h1>Caminos para volver a respirar</h1>
                <p>Recorridos locales para observar, aprender y cuidar.</p>
                <a class="boton" href="#rutas">Ver recorridos</a>
            </div>
        </section>

        <section class="seccion" id="rutas">
            <header class="seccion__titulo">
                <p>Elige tu experiencia</p>
                <h2>Rutas destacadas</h2>
            </header>
            <div class="rejilla-rutas">
                <article class="ruta ruta--principal">
                    <img src="img/bosque.jpg" alt="Sendero entre árboles altos">
                    <div class="ruta__contenido">
                        <p class="nivel">Nivel intermedio</p>
                        <h3>Bosque de niebla</h3>
                        <p>Un trayecto interpretativo entre nacimientos de agua.</p>
                        <a href="#contacto">Consultar salida</a>
                    </div>
                </article>
                <article class="ruta">
                    <img src="img/cascada.jpg" alt="Cascada rodeada de vegetación">
                    <div class="ruta__contenido">
                        <p class="nivel">Nivel básico</p>
                        <h3>Camino del agua</h3>
                        <p>Recorrido breve para reconocer el ciclo hídrico.</p>
                        <a href="#contacto">Consultar salida</a>
                    </div>
                </article>
                <article class="ruta">
                    <img src="img/mirador.jpg" alt="Mirador sobre un valle verde">
                    <div class="ruta__contenido">
                        <p class="nivel">Nivel avanzado</p>
                        <h3>Mirador del valle</h3>
                        <p>Ascenso exigente con observación del paisaje regional.</p>
                        <a href="#contacto">Consultar salida</a>
                    </div>
                </article>
            </div>
        </section>

        <section class="seccion recomendaciones" id="recomendaciones">
            <div>
                <p class="etiqueta">Antes de salir</p>
                <h2>Preparación responsable</h2>
            </div>
            <ul>
                <li>Lleva hidratación y protección solar.</li>
                <li>Regresa todos tus residuos.</li>
                <li>Respeta senderos y fauna local.</li>
            </ul>
        </section>
    </main>
    <footer id="contacto"><p>Ruta Verde · Proyecto académico</p></footer>
</body>
</html>
\end{htmlfile}

\section{Diseño responsivo}

Un diseño responsivo no es una versión pequeña del escritorio. Es una composición
capaz de conservar contenido, lectura y operación en diferentes condiciones de
ancho, orientación, zoom y preferencias del usuario.

\subsection{Viewport y medidas fluidas}

La etiqueta \texttt{meta viewport} permite que el ancho lógico del documento
coincida con el ancho del dispositivo. Ya está incluida en el HTML base.

\begin{csscodedark}
.seccion {
    width: min(100% - 2rem, 1120px);
    margin-inline: auto;
}

.hero h1 {
    font-size: clamp(2.2rem, 7vw, 5rem);
}
\end{csscodedark}

\Idea[Responsivo por naturaleza]{
Antes de escribir una media query, permita que el contenido fluya: use anchos
máximos, porcentajes, \texttt{min()}, \texttt{max()}, \texttt{clamp()}, Flexbox con
\texttt{wrap} y Grid con pistas flexibles.
}

\ExampleBlock[Decisión basada en el contenido]{
Si dos columnas de 300 px, su separación de 24 px y los márgenes laterales ya no
caben con comodidad, ese comportamiento ofrece una razón concreta para cambiar a
una columna. El punto de ruptura nace de la composición, no del nombre de un
dispositivo.
}

\subsection{Imágenes adaptables y \texttt{object-fit}}

\begin{csscodedark}
img {
    display: block;
    max-width: 100%;
}

.ruta img {
    width: 100%;
    height: 220px;
    object-fit: cover;
}
\end{csscodedark}

\begin{center}
\begin{tikzpicture}[
    marco/.style={draw=BrandBlue, very thick, fill=gray!5},
    imagen/.style={draw=PracticeTitle, thick, fill=PracticeBody}
]
    \node[font=\bfseries] at (2.6,3.5) {\texttt{cover}};
    \draw[marco] (0,0.5) rectangle (5.2,3);
    \draw[imagen] (0,0.5) rectangle (5.2,3);
    \draw[dashed, thick] (0.8,0.5) rectangle (4.4,3);
    \node[align=center] at (2.6,1.75) {cubre el marco\\y puede recortar};

    \node[font=\bfseries] at (8.7,3.5) {\texttt{contain}};
    \draw[marco] (6.0,0.5) rectangle (11.2,3);
    \draw[imagen] (6.8,0.5) rectangle (10.4,3);
    \node[align=center] at (8.6,1.75) {muestra todo\\y puede dejar espacio};
\end{tikzpicture}
\end{center}

\subsection{Media queries y mobile first}

En mobile first las reglas base resuelven el ancho pequeño. Las consultas con
\texttt{min-width} añaden cambios cuando existe espacio suficiente. Los puntos de
ruptura deben responder al contenido, no a una lista rígida de dispositivos.

\begin{center}
\begin{tikzpicture}[
    etapa/.style={draw=BrandBlue, thick, rounded corners, minimum width=3cm,
        minimum height=1.1cm, align=center},
    flecha/.style={->, >=Latex, very thick}
]
    \node[etapa, fill=IdeaBody] (m) at (1.8,2) {Reglas base\\una columna};
    \node[etapa, fill=PracticeBody] (t) at (6.2,2)
        {\texttt{min-width: 700px}\\dos columnas};
    \node[etapa, fill=ResultBody] (d) at (10.6,2)
        {\texttt{min-width: 1000px}\\composición amplia};
    \draw[flecha] (m.east) -- (t.west);
    \draw[flecha] (t.east) -- (d.west);
\end{tikzpicture}
\end{center}

\begin{csscodedark}
/* Base: pantalla estrecha */
.encabezado { padding: 1rem; }
.menu { display: flex; flex-wrap: wrap; gap: 0.75rem; }
.recomendaciones { display: grid; gap: 1.5rem; }

@media (min-width: 700px) {
    .recomendaciones {
        grid-template-columns: 1fr 1fr;
        align-items: center;
    }
}

@media (min-width: 1000px) {
    .ruta--principal {
        grid-column: span 2;
    }
}
\end{csscodedark}

\Practice[Ejercicio 2: encontrar el punto de ruptura]{
No elija primero 700 px. Reduzca lentamente el navegador y observe cuándo el
contenido deja de respirar o el menú se vuelve difícil de usar. Registre ese ancho,
ajuste la consulta y explique qué problema concreto resuelve.
}

\Warning[No oculte contenido esencial]{
Una media query puede reorganizar, reducir espacios o cambiar la presentación,
pero no debería eliminar información o acciones necesarias solo para hacer caber
la interfaz.
}

\section{Estados y pseudoelementos}

La interacción debe comunicar qué elemento puede usarse y qué está ocurriendo.
Los estados más habituales son \texttt{:hover}, \texttt{:focus-visible},
\texttt{:active} y \texttt{:focus-within}. Ratón y teclado deben recibir señales
claras.

\begin{csscodedark}
.boton:hover {
    background-color: #0a6048;
}
.boton:focus-visible {
    outline: 3px solid #f7b32b;
    outline-offset: 4px;
}
.boton:active {
    transform: translateY(1px);
}
\end{csscodedark}

Los pseudoelementos \texttt{::before} y \texttt{::after} crean una caja decorativa
asociada a un elemento. No deben contener información esencial que falte en HTML.

\begin{csscodedark}
.nivel::before {
    content: "";
    display: inline-block;
    width: 0.6rem;
    height: 0.6rem;
    margin-right: 0.45rem;
    background: #22a06b;
    border-radius: 50%;
}
\end{csscodedark}


\begin{center}
\begin{tikzpicture}[
    node distance=1.3cm,
    estado/.style={draw=BrandBlue, thick, rounded corners, minimum width=3cm,
        minimum height=0.8cm, fill=IdeaBody, align=center},
    flecha/.style={-{Latex[length=3mm]}, thick}
]
    \node[estado] (normal) {estado normal};
    \node[estado, right=of normal] (hover) {\texttt{:hover}};
    \node[estado, right=of hover] (focus) {\texttt{:focus}};
    \draw[flecha] (normal) -- node[above]{puntero} (hover);
    \draw[flecha] (normal) to[bend right=22] node[below]{teclado} (focus);
\end{tikzpicture}
\end{center}

\section{Transformaciones y transiciones}

\texttt{transform} modifica la representación visual sin reorganizar el flujo.
Puede trasladar, escalar, rotar o inclinar. \texttt{transition} interpola el cambio
entre un estado inicial y uno final.

\begin{csscodedark}
.ruta {
    transition:
        transform 220ms ease,
        box-shadow 220ms ease;
}

.ruta:hover {
    transform: translateY(-6px);
    box-shadow: 0 18px 40px #14362a28;
}
\end{csscodedark}

\begin{center}
\begin{tikzpicture}[
    estado/.style={draw=BrandBlue, thick, rounded corners, minimum width=3.4cm,
        minimum height=1.2cm, align=center},
    flecha/.style={<->, >=Latex, very thick}
]
    \node[estado, fill=IdeaBody] (a) at (0,0) {estado inicial\\\texttt{translateY(0)}};
    \node[estado, fill=ResultBody] (b) at (6,0) {estado final\\\texttt{translateY(-6px)}};
    \draw[flecha] (a) -- node[above]{220 ms · \texttt{ease}} (b);
\end{tikzpicture}
\end{center}

\DefBlock[Cuatro partes de una transición]{
La propiedad que cambia, la duración, la función de tiempo y el retraso opcional.
Evite \texttt{transition: all}; declarar propiedades concretas hace el efecto más
predecible y evita animar cambios costosos o accidentales.
}

\subsection{Transformaciones habituales}

\begin{tabularx}{\linewidth}{>{\ttfamily\bfseries\small}p{4.4cm}X}
\toprule
Función & Resultado \\
\midrule
translate(x, y) & Desplaza visualmente. \\
scale(n) & Amplía o reduce desde el origen de transformación. \\
rotate(ángulo) & Gira el elemento. \\
skew(ángulo) & Inclina sus ejes. \\
transform-origin & Cambia el punto desde el que ocurre la transformación. \\
\bottomrule
\end{tabularx}

\Practice[Ejercicio 3: microinteracción]{
Cree un botón con tres estados: reposo, foco visible y activación. Use una sola
transformación pequeña y una transición entre 150 y 300 ms. Compruebe que el
texto no se mueve y que el foco no depende únicamente del color.
}

\section{Animaciones con \texttt{@keyframes}}

Una animación describe uno o más estados intermedios y puede ejecutarse sin que
una pseudoclase mantenga el cambio. Es adecuada cuando hay una secuencia temporal;
para un cambio simple entre dos estados suele bastar una transición.

\begin{csscodedark}
@keyframes aparecer {
    from {
        opacity: 0;
        transform: translateY(18px);
    }
    to {
        opacity: 1;
        transform: translateY(0);
    }
}

.ruta {
    animation: aparecer 500ms ease-out both;
}
.ruta:nth-child(2) { animation-delay: 100ms; }
.ruta:nth-child(3) { animation-delay: 200ms; }
\end{csscodedark}

\begin{center}
\begin{tikzpicture}[
    estado/.style={draw=BrandBlue, thick, rounded corners, minimum width=2.5cm,
        minimum height=1cm, align=center},
    flecha/.style={->, >=Latex, thick}
]
    \node[estado, fill=WarnBody] (a) at (0,0) {0\%\\oculto y abajo};
    \node[estado, fill=IdeaBody] (b) at (4.3,0) {50\%\\en movimiento};
    \node[estado, fill=ResultBody] (c) at (8.6,0) {100\%\\visible y ubicado};
    \draw[flecha] (a.east) -- (b.west);
    \draw[flecha] (b.east) -- (c.west);
\end{tikzpicture}
\end{center}

\begin{tabularx}{\linewidth}{>{\ttfamily\bfseries\small}p{6.3cm}X}
\toprule
Propiedad & Control \\
\midrule
animation-name & Nombre de los \texttt{@keyframes}. \\
animation-duration & Tiempo de un ciclo. \\
animation-timing-function & Ritmo del cambio. \\
animation-delay & Espera antes de comenzar. \\
animation-iteration-count & Número de repeticiones. \\
animation-direction & Sentido de cada ciclo. \\
animation-fill-mode & Estado aplicado antes o después. \\
animation-play-state & Ejecución o pausa. \\
\bottomrule
\end{tabularx}

\Compare[Transición o animación]{
Use transición cuando existe un estado inicial y otro final provocado por una
interacción. Use animación cuando necesita una secuencia, repetición, retrasos
independientes o varios puntos intermedios.
}

\section{Efectos visuales con moderación}

Sombras, degradados, opacidad y filtros pueden mejorar profundidad o legibilidad,
pero también reducir contraste y rendimiento si se acumulan.

\begin{csscodedark}
.hero {
    color: white;
    background:
        linear-gradient(90deg, #073b2ddd, #073b2d44),
        url("../img/portada.jpg") center / cover;
}

.ruta img {
    filter: saturate(0.9) contrast(1.05);
    transition: filter 220ms ease, transform 220ms ease;
}
.ruta:hover img {
    filter: saturate(1.1) contrast(1.08);
    transform: scale(1.025);
}
\end{csscodedark}

\Warning[Opacidad y texto]{
Aplicar \texttt{opacity} al contenedor vuelve transparentes también sus hijos. Para
oscurecer una imagen sin afectar el texto, use una capa, un pseudoelemento o un
degradado en el fondo.
}

\section{Movimiento accesible}

No todas las personas toleran o desean movimiento. La consulta
\texttt{prefers-reduced-motion} permite reducirlo cuando el sistema expresa esa
preferencia.

\begin{csscodedark}
@media (prefers-reduced-motion: reduce) {
    html {
        scroll-behavior: auto;
    }

    *, *::before, *::after {
        animation-duration: 0.01ms !important;
        animation-iteration-count: 1 !important;
        transition-duration: 0.01ms !important;
        scroll-behavior: auto !important;
    }
}
\end{csscodedark}

\Idea[Movimiento con propósito]{
Una animación útil orienta la atención, confirma una acción o explica una
transición espacial. Si no cumple una función reconocible, probablemente puede
eliminarse.
}

\section{CSS integrado del laboratorio}

La siguiente síntesis muestra cómo conviven Grid, Flexbox, reglas responsivas y
movimiento. Las reglas tipográficas menores pueden ampliarse durante la práctica.

\begin{csscodedark}
:root { --verde: #0a6048; --oscuro: #102820; --claro: #f2f7f4; --acento: #f7b32b; }
*, *::before, *::after { box-sizing: border-box; }
html { scroll-behavior: smooth; }
body { margin: 0; color: var(--oscuro); background: var(--claro); font-family: Arial, sans-serif; line-height: 1.6; }
img { display: block; max-width: 100%; }
a { color: inherit; }

.encabezado { display: flex; justify-content: space-between; align-items: center; flex-wrap: wrap; gap: 1rem; padding: 1rem max(1rem, calc((100% - 1120px) / 2)); background: white; }
.marca { color: var(--verde); font-weight: 800; text-decoration: none; }
.menu { display: flex; flex-wrap: wrap; gap: 0.75rem 1.25rem; margin: 0; padding: 0; list-style: none; }
.menu a:focus-visible, .boton:focus-visible, .ruta a:focus-visible { outline: 3px solid var(--acento); outline-offset: 4px; }

.hero { min-height: 70vh; display: grid; align-items: center; padding: 4rem max(1rem, calc((100% - 1120px) / 2)); color: white; background: linear-gradient(90deg, #073b2dee, #073b2d44), url("../img/portada.jpg") center / cover; }
.hero__contenido { max-width: 720px; }
.hero h1 { margin: 0; font-size: clamp(2.2rem, 7vw, 5rem); line-height: 1.05; }
.boton { display: inline-block; padding: 0.75rem 1rem; color: white; background: var(--verde); border-radius: 999px; text-decoration: none; transition: background-color 180ms ease, transform 180ms ease; }
.boton:hover { background: #073b2d; transform: translateY(-2px); }

.seccion { width: min(100% - 2rem, 1120px); margin-inline: auto; padding-block: 4rem; }
.seccion__titulo { margin-bottom: 2rem; }
.rejilla-rutas { display: grid; grid-template-columns: repeat(auto-fit, minmax(min(100%, 260px), 1fr)); gap: 1.5rem; }
.ruta { overflow: hidden; display: flex; flex-direction: column; background: white; border-radius: 1rem; box-shadow: 0 8px 24px #14362a18; animation: aparecer 500ms ease-out both; transition: transform 220ms ease, box-shadow 220ms ease; }
.ruta:nth-child(2) { animation-delay: 100ms; }
.ruta:nth-child(3) { animation-delay: 200ms; }
.ruta:hover { transform: translateY(-6px); box-shadow: 0 18px 40px #14362a28; }
.ruta img { width: 100%; height: 220px; object-fit: cover; }
.ruta__contenido { flex: 1; display: flex; flex-direction: column; padding: 1.25rem; }
.ruta__contenido a { margin-top: auto; color: var(--verde); font-weight: 700; }
.nivel::before { content: ""; display: inline-block; width: 0.6rem; height: 0.6rem; margin-right: 0.45rem; background: #22a06b; border-radius: 50%; }
.recomendaciones { display: grid; gap: 1.5rem; }
footer { padding: 2rem; color: white; background: var(--oscuro); text-align: center; }

@keyframes aparecer {
    from { opacity: 0; transform: translateY(18px); }
    to { opacity: 1; transform: translateY(0); }
}
@media (min-width: 700px) {
    .recomendaciones { grid-template-columns: 1fr 1fr; align-items: center; }
}
@media (min-width: 1000px) {
    .ruta--principal { grid-column: span 2; }
}
@media (prefers-reduced-motion: reduce) {
    html { scroll-behavior: auto; }
    *, *::before, *::after { animation-duration: 0.01ms !important; animation-iteration-count: 1 !important; transition-duration: 0.01ms !important; }
}
\end{csscodedark}

\section{Práctica guiada: agenda cultural}

Construya una cartelera de actividades culturales con un evento destacado y al
menos cinco eventos secundarios.

\begin{enumerate}
    \item Diseñe primero una versión de una columna, completamente legible.
    \item Cree una cuadrícula fluida con \texttt{auto-fit} y \texttt{minmax()}.
    \item En ancho amplio, haga que el evento destacado abarque dos columnas.
    \item Organice el contenido interno de cada tarjeta con Flexbox.
    \item Añada una consulta basada en el punto donde el contenido necesita cambiar.
    \item Cree estados de foco y ratón equivalentes.
    \item Incorpore una transición breve y una animación de entrada no repetitiva.
    \item Implemente la preferencia de movimiento reducido.
\end{enumerate}

\Practice[Prueba en cuatro condiciones]{
Revise la interfaz a 360, 700, 1000 y 1400 px. En cada ancho compruebe lectura,
orden, navegación, desbordamiento, tamaño de imágenes, foco por teclado y reducción
de movimiento.
}

\section{Ejercicio autónomo: destinos regionales}

Desarrolle una página para presentar mínimo seis destinos regionales. Incluya una
portada, categorías visuales, galería de destinos, recomendaciones y pie de página.

\subsection{Requisitos mínimos}

\begin{itemize}
    \item HTML semántico y contenido completo antes de aplicar efectos.
    \item Grid para la galería y Flexbox para al menos dos componentes internos.
    \item Uso de \texttt{fr}, \texttt{repeat()}, \texttt{minmax()} y separación con \texttt{gap}.
    \item Un elemento que abarque más de una columna en ancho amplio.
    \item Imágenes fluidas con una decisión justificada entre \texttt{cover} y \texttt{contain}.
    \item Enfoque mobile first y mínimo una media query basada en el contenido.
    \item Estados visibles para ratón y teclado.
    \item Una transición y una animación breve con propósito distinto.
    \item Compatibilidad con \texttt{prefers-reduced-motion}.
\end{itemize}

\Warning[Los adicionales no sustituyen la base]{
Agregar más destinos, animaciones, filtros o efectos no exime del cumplimiento de
los requisitos mínimos. Complete y compruebe primero lo solicitado; después añada
los elementos complementarios sin deteriorar accesibilidad ni rendimiento.
}

\section{Preguntas de comprensión}

\begin{enumerate}
    \item ¿Qué diferencia conceptual existe entre Flexbox y Grid?
    \item ¿Qué representa una unidad \texttt{fr}?
    \item ¿Cuándo crea Grid filas implícitas?
    \item ¿Qué diferencia existe entre alinear elementos y distribuir la cuadrícula?
    \item ¿Cómo se comportan \texttt{auto-fit} y \texttt{minmax()} juntos?
    \item ¿Por qué un punto de ruptura debe responder al contenido?
    \item ¿Qué diferencia visual existe entre \texttt{cover} y \texttt{contain}?
    \item ¿Cuándo conviene una transición y cuándo una animación?
    \item ¿Qué propósito cumple \texttt{prefers-reduced-motion}?
\end{enumerate}

\section{Errores frecuentes y diagnóstico}

\begin{tabularx}{\linewidth}{X X}
\toprule
Síntoma & Revisión recomendada \\
\midrule
Las columnas desbordan. & Revise mínimos, contenido largo y \texttt{minmax(0, 1fr)}. \\
Un área no aparece. & Compruebe nombres y que el área sea rectangular. \\
\texttt{align-content} no cambia nada. & Verifique que la cuadrícula sea menor que su contenedor. \\
La media query nunca se activa. & Revise sintaxis, unidad y orden de reglas. \\
La imagen se deforma. & Use dimensiones coherentes y \texttt{object-fit}. \\
La transición es brusca. & Defina estado inicial, propiedad concreta y duración. \\
La animación vuelve al inicio. & Revise \texttt{animation-fill-mode}. \\
El movimiento distrae. & Reduzca distancia, duración y repeticiones; respete preferencias. \\
\bottomrule
\end{tabularx}

\section{Lista de verificación}

\begin{itemize}
    \ChecklistItem{La página funciona en una columna antes de las media queries.}
    \ChecklistItem{La cuadrícula tiene columnas y separaciones comprensibles.}
    \ChecklistItem{Las líneas o áreas se usan sin alterar el orden lógico del HTML.}
    \ChecklistItem{No existe desplazamiento horizontal accidental.}
    \ChecklistItem{Las imágenes conservan proporción y encuadre apropiado.}
    \ChecklistItem{Los puntos de ruptura resuelven problemas observados.}
    \ChecklistItem{La navegación funciona con teclado y el foco es visible.}
    \ChecklistItem{Las transiciones declaran propiedades concretas.}
    \ChecklistItem{Las animaciones tienen propósito y duración moderada.}
    \ChecklistItem{Se implementó movimiento reducido.}
    \ChecklistItem{Los requisitos mínimos se fijaron antes de los adicionales.}
\end{itemize}

\section{Consulta rápida}

\begin{multicols}{2}
\textbf{Grid y responsive}
\begin{itemize}
    \item \texttt{grid-template-columns}
    \item \texttt{grid-template-rows}
    \item \texttt{gap}, \texttt{fr}, \texttt{repeat()}
    \item \texttt{minmax()}, \texttt{auto-fit}
    \item \texttt{grid-column}, \texttt{grid-area}
    \item \texttt{place-items}, \texttt{place-self}
    \item \texttt{@media (min-width: ...)}
    \item \texttt{clamp()}, \texttt{object-fit}
\end{itemize}

\columnbreak
\textbf{Interacción y movimiento}
\begin{itemize}
    \item \texttt{:hover}, \texttt{:focus-visible}, \texttt{:active}
    \item \texttt{::before}, \texttt{::after}
    \item \texttt{transform}, \texttt{transform-origin}
    \item \texttt{transition-property}, \texttt{duration}, \texttt{timing-function}
    \item \texttt{@keyframes}, \texttt{animation}
    \item \texttt{filter}, \texttt{opacity}, sombras y degradados
    \item \texttt{prefers-reduced-motion}
\end{itemize}
\end{multicols}


\section{Cierre de la guía}
\Result[Cierre]{
El propósito de esta guía es que el aprendiz pueda regresar a ella durante la práctica sin encontrar todavía la solución de los temas de la guía siguiente. Antes de avanzar, debe ser capaz de explicar con sus propias palabras los conceptos centrales y construir una solución funcional sin copiar el ejemplo completo.
}
\end{document}
